\chapter{Background}

% Related work from ITSC

We have been primarily motivated
to study games of ordered preference
by Zanardi et al. \cite{zanardi:2022} and Lee et al. \cite{lee:2025}.
We explicitly build on this research
to compute for a receding horizon,
rather than a finite-time horizon,
to improve the scalability of these games.
We have additionally found insights in receding horizon games
\cite{hall:2022,hall:2024}, although the formulation is not directly applicable
to the problem of preference ordering.

Prior work, closely related to the formulation
of games of ordered preference in this paper,
is on \emph{prioritized metrics} 
and \emph{game-theoretic planning}.

\smallskip \noindent
\textsc{prioritized metrics} \quad 
Posetal games \cite{zanardi:2022} formalize games
with prioritized metrics
through partial orders over system designs \cite{censi:2019}, using symbolic structures to encode constraint and preference hierarchies \cite{wongpiromsarn:2021}.
While posetal games focus
on discrete decision spaces, we adopt lexicographic optimization \cite{penlington:2024},
that extends prioritized metrics to continuous action spaces \cite{lee:2025}.
However, the lexicographic nested problem structure introduces computational challenges: lower-priority objectives are optimized only after higher-priority ones are satisfied, creating a hierarchy of interdependent subproblems. To mitigate this computational intractability, prior work transcribes the lexicographic hierarchy into mathematical programs with complementarity constraints \cite{nurkanovic:2024}, with the addition of relaxation schemes \cite{lee:2025}.
Yet, coupled constraints cause scalability to remain limited.
By framing the problem
within a receding horizon setting, we exploit the temporal decoupling inherent
to IBR algorithms, which are known to be computational tractable.

\smallskip \noindent
\textsc{game-theoretic planning}\quad Game-theoretic traffic
management is needed to analyze the interactive nature of preferences between
autonomous agents \cite{qin:2024}. Efficient solution strategies exist
for game-theoretic planning \cite{laine:2021,williams:2023,mehr:2023,zhu:2023}, but solving games that involve nested optimization problems typically requires approximate solutions, similar to those used in solving trajectory games with continuous action spaces \cite{jond:2022,peters:2022,liu:2023}.
Moreover, predictions of other agents' intentions combined with planning
increase the accuracy of approximate solutions \cite{li:2022,zhang:2023}. Receding
horizon in game-theoretic planning has been studied in the classical formulation
\cite{veer:2023} and in variants---where objects are divided into immediate
and mid-term metrics \cite{fisac:2019}. These approaches do not consider the
problem of preferences in agent interactions.
